%%%%%%%%%%%%%%%%%%%%%%%%%%%%%%%%%%%%%%%%%%%%%%%%%%%%%%%%%%%%%%%%%%%%%%%%
%% Unrestricted Chow-Liu Structure Selection for Mixed
%% Discrete-Continuous Data
%% Final manuscript for Behaviormetrika
%%
%% This file compiles standalone with plain article class.
%% To submit to Behaviormetrika, replace the \documentclass line below
%% with the Springer class (\documentclass{sn-jnl} or similar) per the
%% journal's author kit. The body content (after \begin{document}) is
%% portable.
%%%%%%%%%%%%%%%%%%%%%%%%%%%%%%%%%%%%%%%%%%%%%%%%%%%%%%%%%%%%%%%%%%%%%%%%

\documentclass[10pt,a4paper]{article}

\usepackage[margin=1in]{geometry}
\usepackage{amsmath,amssymb,amsthm}
\usepackage{booktabs}
\usepackage{graphicx}
\graphicspath{{./}}  % all figures sit next to manuscript.tex (flat, no subfolder)
% Numbered references in square brackets, sorted and range-compressed,
% so \citep{a,b,c} renders as [1, 2, 3] (or [1-3] when consecutive).
\usepackage[numbers,square,sort&compress]{natbib}
\usepackage{url}

% Make citation numbers and cross-references clickable. Links are
% rendered in dark colours (academic-paper style): citations in dark
% green, internal cross-references (Proposition, Corollary, Section,
% Figure, Table) in dark blue, URLs in muted blue. The colours are
% subtle enough not to distract but visible enough for the reader to
% see which text is a click target.
\usepackage{hyperref}
\hypersetup{
  colorlinks   = true,
  linkcolor    = {[rgb]{0.0, 0.0, 0.55}},
  citecolor    = {[rgb]{0.0, 0.45, 0.0}},
  urlcolor     = {[rgb]{0.0, 0.0, 0.7}},
  bookmarksopen = true,
  bookmarksnumbered = true
}

% cleveref must be loaded AFTER hyperref. It makes references like
% \cref{prop:main} render as "Proposition 1" with the whole phrase
% being a clickable link, rather than "Proposition 1" with only the
% number being the link. The capitalize option keeps the leading word
% capitalised (Proposition, Corollary, Section, ...) at sentence start.
\usepackage[capitalize,noabbrev]{cleveref}
\crefname{section}{Section}{Sections}
\crefname{subsection}{Section}{Sections}
\crefname{table}{Table}{Tables}
\crefname{figure}{Figure}{Figures}
\crefname{equation}{equation}{equations}

% Algorithm box without external package (portable).
% If you have algorithm/algorithmic packages available, you can replace
% the \begin{algBox}...\end{algBox} environment below with a standard
% \begin{algorithm}...\end{algorithm} block.
\newenvironment{algBox}[1]
  {\medskip\par\noindent\textbf{Algorithm 1: #1}\\\rule{\textwidth}{0.5pt}\par\small}
  {\par\rule{\textwidth}{0.5pt}\medskip}

% Theorem environments
\newtheorem{proposition}{Proposition}
\newtheorem{corollary}{Corollary}
\theoremstyle{remark}
\newtheorem*{remark}{Remark}

%%%%%%%%%%%%%%%%%%%%%%%%%%%%%%%%%%%%%%%%%%%%%%%%%%%%%%%%%%%%%%%%%%%%%%%%

\title{A Note on Unrestricted Chow--Liu Tree Selection for Mixed
Discrete--Continuous Data}

\author{Anonymous}

\date{}

\begin{document}
\maketitle

%%%%%%%%%%%%%%%%%%%%%%%%%%%%%%%%%%%%%%%%%%%%%%%%%%%%%%%%%%%%%%%%%%%%%%%%
\begin{abstract}
The Chow--Liu algorithm summarises the dependency structure of
mixed discrete--continuous data as a tree of strongest pairwise
associations. Earlier versions exclude trees in which a
continuous variable connects two discrete variables, because such
a path complicates the later density-fitting step. We show that
this restriction is unnecessary for tree selection, which depends
only on pairwise mutual-information values, well-defined for
every combination of variable types. The restriction applies only
to the later density-fitting step. We prove consistency of the unrestricted
procedure and give finite-sample error bounds under a tail
assumption on the pairwise mutual-information estimates. On synthetic mixed graphs and a
synthetic gene-network benchmark, the unrestricted procedure
recovers trees that the restricted procedure cannot. On the
fully discrete Asia Bayesian benchmark, where Edwards'
restriction is inactive by construction, the pipeline attains the
tree-approximation ceiling. A breast-cancer gene-expression
dataset is included as an illustration, since no ground truth is
available. The comparison with a discretisation-based estimator is
regime-dependent. In the sample-size sweep, the discretisation
estimator performs better at $n = 75$, the
$k$-nearest-neighbour estimator performs better at
$n = 1{,}500$, and the two are statistically indistinguishable
at $n \in \{100, 200, 500, 1{,}000\}$.
\end{abstract}

\medskip
\noindent\textbf{Keywords:} Chow--Liu algorithm; mixed graphical
models; mutual information; structure learning; conditional Gaussian
model.

%%%%%%%%%%%%%%%%%%%%%%%%%%%%%%%%%%%%%%%%%%%%%%%%%%%%%%%%%%%%%%%%%%%%%%%%
\section{Introduction}
\label{sec:intro}

Consider a study that records a panel of $p$ variables for $n$
individuals. Some variables are categorical, such as genotypes at
single nucleotide polymorphism (SNP) sites, and some are continuous,
such as gene-expression levels or disease biomarkers. A common first
task is to summarise the pairwise dependencies among these
variables. The Chow--Liu algorithm \citep{chow1968approximating} answers a
tractable version of this task. It returns the tree-structured graph whose edges together capture as
much of the joint dependency as possible, measured by the
sum of pairwise mutual informations (MI). The optimal tree is the maximum-weight spanning tree (MST) of
the complete graph whose edge weights are pairwise MI values.
Equivalently, the chosen tree minimises the
Kullback--Leibler (KL) divergence to the true joint distribution
among all tree-structured approximations.

This mixed-type setting arises naturally in genetics
\citep{rockman2006genetics, nica2013expression} and in
epidemiology where disease status accompanies continuous
biomarkers. The conditional Gaussian (CG) model
\citep{lauritzen1989graphical, lauritzen1996graphical} is the standard
joint distribution for such data. A modelling issue arises in
the CG framework. If the tree contains a path $D_1$-$G$-$D_2$ in which
$G$ is continuous and $D_1, D_2$ are discrete, then the tree
factorisation requires the local conditional $P(D_2 \mid G)$, which
is not part of the standard CG parametrisation
\citep{whittaker1990graphical}. To avoid
this difficulty, \citet{edwards2010selecting} imposed a structural
restriction on the tree search. The restriction forbids continuous
nodes from separating discrete nodes. It guarantees a fully
parametrised CG model on the learned tree, but it also prevents the
algorithm from recovering $D$-$G$-$D$ Markov paths that arise
naturally in mixed-type data. \citet{suzuki2017chow} subsequently evaded the restriction in
practice by discretising continuous
variables with a nested mesh and applying a minimum description
length (MDL) penalty to the discretised mutual information. The
restriction is sometimes treated as part of the mixed-type
Chow--Liu procedure, although it constrains only the later CG
parametrisation step.

We observe that the restriction belongs to the later
density-fitting step, not to the Chow--Liu objective itself. The
objective is a sum of pairwise mutual informations, each of which is
symmetric in its two arguments. The objective therefore does not require any directional factor such
as $P(D \mid G)$. Under
standard regularity conditions, the unrestricted procedure recovers
the population Chow--Liu tree as the sample size grows. We state
this consistency result in \cref{prop:main} and a forest analogue
under a shrinking threshold in \cref{prop:forest}. A finite-sample
bound (\cref{cor:tree}) and a tighter per-swap bound
(\cref{cor:swap}) follow. Our contribution is the clarification, its
empirical validation, and the supporting finite-sample analysis. The
proof techniques are standard.

\paragraph{Contributions.}
\begin{enumerate}
\item \emph{Separation of structure selection from CG
parametrisation.} The Chow--Liu objective is a sum of pairwise mutual
informations. Each is a symmetric quantity that can be written using
either $P(G \mid D)$ or $P(D \mid G)$ through the standard entropy
decomposition. The objective therefore does not require modelling
$P(D \mid G)$ for any pair. Edwards' restriction applies only to
the subsequent step of fitting a fully parametric CG density on
the selected tree.

\item \emph{An estimator-agnostic procedure with finite-sample
bounds.} Algorithm~1 uses the
Kraskov--St\"ogbauer--Grassberger (KSG) $k$-nearest-neighbour estimator
for continuous--continuous pairs \citep{kraskov2004estimating}, the
mixed-pair $k$-NN estimator of \citet{gao2017estimating} for
continuous--discrete pairs, and the empirical plug-in for
discrete--discrete pairs. \Cref{prop:main} establishes consistency
whenever the pairwise mutual-information estimates are per-pair
consistent. Under an additional sub-Gaussian deviation assumption,
\cref{cor:tree,cor:swap} give finite-sample tree-recovery bounds.
The per-swap bound in \cref{cor:swap} uses the fact that exchanging
one edge for another in a maximum spanning tree creates exactly one
cycle, and bounds the probability of a single incorrect exchange. It
improves the coarser tree-symmetric-difference bound in
\cref{cor:tree} by a factor of $(p-1)^2$ in the exponent.

\item \emph{Regime-dependent empirical comparison.} We compare
the proposed pipeline (with the KSG MI estimator as input)
against the Edwards-restricted MST baseline and against the
MDL-discretisation MI estimator of \citet{suzuki2017chow} over a
range of graph sizes, effect sizes, conditional families, and
sample sizes. The proposed pipeline outperforms the Edwards-restricted
baseline across the tested graph-size and effect-size settings on
graphs containing $D$-$G$-$D$ paths. The comparison with the
MDL-discretisation estimator is sample-size dependent. MDL
discretisation performs better at $n = 75$ (where its bin selection
acts as a form of model selection), the two estimators are
statistically indistinguishable at $n \in \{100, 200, 500, 1{,}000\}$,
and KSG performs better at $n = 1{,}500$.
\end{enumerate}

Our results cover only the first step of the Chow--Liu workflow,
selecting which edges form the tree, and the returned graph is
undirected. A learned edge means that two variables are
statistically associated; it carries no causal interpretation. A
separate step, when needed, writes the joint distribution as an
explicit parametric density on the chosen tree. The standard
conditional-Gaussian density uses local factors of the form
$P(G \mid D)$, but not $P(D \mid G)$. If the chosen tree contains a
$D$-$G$-$D$ path, this density-fitting step needs the factor
$P(D \mid G)$, which lies outside the standard parametrisation. That step still requires either Edwards' restriction, which
forbids the path in the first place, or an alternative
parametrisation of $P(D \mid G)$. We do not address this
density-fitting step.

%%%%%%%%%%%%%%%%%%%%%%%%%%%%%%%%%%%%%%%%%%%%%%%%%%%%%%%%%%%%%%%%%%%%%%%%
\section{Background}
\label{sec:background}

\subsection{Chow--Liu algorithm}
\label{subsec:chowliu}

A graphical model represents the joint distribution of $p$ random
variables $X_1, \ldots, X_p$ as a graph whose edges encode pairwise
dependencies. When the graph is constrained to be a tree, each
variable is connected to the rest of the graph through a single path,
and the joint distribution factorises into a product of pairwise
marginals along the tree edges and univariate marginals at the
nodes. The Chow--Liu algorithm seeks the tree $T^*$ that best
approximates the joint distribution $P$ in Kullback--Leibler (KL)
divergence:
\begin{equation}
T^* \;=\; \arg\max_{T \in \mathcal{T}} \sum_{(i,j) \in T} I(X_i; X_j),
\label{eq:chowliu_obj}
\end{equation}
where $\mathcal{T}$ is the set of all spanning trees on $p$ vertices
and $I(X_i; X_j)$ is the pairwise mutual information between $X_i$
and $X_j$. The right-hand side of (\ref{eq:chowliu_obj}) is the
total information captured by a tree, so the procedure selects the
tree that captures as much pairwise dependency as possible. The
optimisation reduces to Kruskal's algorithm
\citep{kruskal1956shortest} applied to the complete graph weighted
by the pairwise mutual informations. Kruskal's algorithm adds edges
in decreasing order of weight, skipping any edge that would form a
cycle, until $p - 1$ edges have been added. The classical Chow--Liu algorithm therefore
returns a spanning tree even when some pairwise mutual informations
are estimated as zero or near-zero. In the thresholded variant used
below (\cref{prop:forest}), edges with estimated weight at or below
a threshold $\tau_n$ are excluded before the Kruskal step is run, so
the output may be a forest rather than a connected tree.
Tree-structured graphical models remain a standard component of the
probabilistic graphical models toolkit
\citep{koller2009probabilistic} because of this clean optimality
property and the efficiency of Kruskal's algorithm.

\subsection{Mixed-type mutual information estimation}
\label{subsec:mi}

Mutual information is well-defined for any pair of random variables,
but its estimation method depends on whether the variables are
continuous, discrete, or mixed. For purely continuous pairs under a
joint Gaussian assumption, the mutual information has the closed form
$I(X; Y) = -\tfrac{1}{2}\log(1 - \rho^2)$, where $\rho$ is the
Pearson correlation coefficient \citep{cover2006elements}. For purely
discrete pairs, the mutual information is estimated by the empirical
plug-in, in which the joint and marginal probabilities are replaced
by their sample frequencies. For mixed pairs, where one variable is
continuous and the other discrete, neither of these forms applies. A
distribution-free alternative is the $k$-nearest-neighbour estimator of
\citet{kraskov2004estimating}, which builds on the
\citet{kozachenko1987sample} entropy estimator and was extended to
mixed discrete--continuous variables with provably consistent
convergence rates by \citet{gao2017estimating}. The method uses
distances between each sample point and its $k$-th nearest neighbour,
and is consistent under standard regularity conditions on the
underlying distribution. We use $k = 5$ as the default value
throughout. For continuous--continuous pairs in which a joint Gaussian
assumption is not appropriate, we use the Kraskov--St\"ogbauer--%
Grassberger (KSG) estimator \citep{kraskov2004estimating}, a
$k$-nearest-neighbour estimator from the same paper that is designed
for the two-continuous-variable case.

\subsection{Edwards' restriction}
\label{subsec:edwards}

In the conditional Gaussian model, the tree factorisation takes the
form $P(X_1, \ldots, X_p) = \prod_i P(X_i \mid X_{\mathrm{pa}(i)})$,
where $\mathrm{pa}(i)$ denotes the parent of node $i$ in the tree
rooted at any node. When each node is conditioned on a parent of the same type, the
conditional distributions on the right-hand side are standard. The
difficulty arises when the parent and child have different types. For a tree containing a
path $D_1$-$G$-$D_2$ rooted at $D_1$, the factorisation becomes
$P = P(D_1) \cdot P(G \mid D_1) \cdot P(D_2 \mid G)$. The first two
factors are standard in the conditional Gaussian model. $P(D_1)$ is
a discrete distribution and $P(G \mid D_1)$ is the Gaussian
conditional that defines the model. The third factor,
$P(D_2 \mid G)$, is the conditional distribution of a discrete
variable given a continuous one, which is not part of the standard
CG parametrisation. \citet{edwards2010selecting} avoided this
difficulty by restricting Kruskal's algorithm during tree
selection. The modified Kruskal's algorithm rejects any candidate
edge whose addition would place a continuous node between two
discrete components. The resulting tree always admits a full
parametrisation within the standard CG model, but the procedure
cannot recover any $D$-$G$-$D$ Markov chain, even when the
data-generating distribution contains one. \Cref{fig:problem} illustrates the
configuration that the restriction forbids.

\begin{figure}[ht]
\centering
\includegraphics[width=0.85\textwidth]{Fig01_problem.png}
\caption{The $D$-$G$-$D$ configuration that Edwards' restriction
forbids by construction. A continuous variable $G$ lies on the
path between two discrete variables $D_1$ and $D_2$, so the
factorisation requires $P(D_2 \mid G)$, which is not part of the
standard CG parametrisation.}
\label{fig:problem}
\end{figure}

\subsection{Related work}
\label{subsec:related}

Structure learning for graphical models is a well-developed area,
with methods ranging from constraint-based to score-based and from
parametric to nonparametric. \citet{drton2017structure} provide a
recent overview of the field. Within mixed-type graphical modelling,
\citet{yang2014mixed} proposed $\ell_1$-regularised models in which
the joint distribution is represented as a member of an exponential
family. Their approach handles arbitrary graph structures rather than
being restricted to trees, but it targets a parametric Markov
random field
and requires convex optimisation rather than a combinatorial tree
search. On the mutual information side,
\citet{belghazi2018mine} and \citet{poole2019variational} developed
neural variational lower bounds on mutual information (the
mutual-information neural estimator, MINE, and the
noise-contrastive estimator, InfoNCE) used in representation
learning. \citet{mcallester2020formal} showed that any such
variational bound has fundamental limitations in the
distribution-free setting. Matching the true mutual information
requires sample sizes that grow exponentially in the mutual
information itself. The closest prior work to ours is
\citet{suzuki2017chow}, who proposed an MDL-discretisation MI
estimator and used it inside an unrestricted Chow--Liu procedure,
thereby lifting Edwards' restriction in practice via discretisation
of the continuous variables. We discuss the MDL-discretisation
procedure in detail and benchmark against it empirically in
\cref{subsec:suzuki}.

%%%%%%%%%%%%%%%%%%%%%%%%%%%%%%%%%%%%%%%%%%%%%%%%%%%%%%%%%%%%%%%%%%%%%%%%
\section{Theoretical statements}
\label{sec:theory}

\subsection{Consistency of the unrestricted Chow--Liu procedure}
\label{subsec:consistency}

We refer to two quantities throughout. The first is the unrestricted
Chow--Liu tree, obtained by maximising the total pairwise mutual
information over all spanning trees. \Cref{prop:main} concerns this
quantity. The second is a thresholded maximum spanning forest,
obtained by first discarding edges with weight at or below a
threshold $\tau_n$ and then applying Kruskal's algorithm to the
remaining edges. \Cref{prop:forest} concerns this
thresholded forest under additional separation assumptions. The
zero-threshold rule used in many of our experiments returns a forest
when some empirical mutual information estimates are zero or
negative, but it is not a direct instance of
\cref{prop:forest}.

\begin{proposition}[Unrestricted Chow--Liu is well-defined for mixed
types]
\label{prop:main}
Let $X_1, \ldots, X_p$ be random variables of arbitrary types. Assume

\begin{itemize}
\item[(A1)] the pairwise mutual informations $I(X_i; X_j)$ are
well-defined and finite for all pairs;
\item[(A2)] the maximum-weight spanning tree is unique;
\item[(A3)] for each pair, $\hat{I}_n(X_i; X_j) \overset{p}{\to}
I(X_i; X_j)$ as $n \to \infty$.
\end{itemize}

Then the unrestricted Chow--Liu algorithm recovers $T^*$ with
probability tending to $1$, regardless of whether $T^*$ contains
$D$-$G$-$D$ paths.
\end{proposition}

The proof follows the standard M-estimation argument over the finite
set of spanning trees \citep{vandervaart2000asymptotic}. Uniform
convergence over $\mathcal{T}$ follows from (A3) applied to each of
the $\binom{p}{2}$ pairs, and identification follows from (A2). The
essential ingredient is the symmetry identity in
\cref{eq:identity}, recorded as a remark after the proof, which
makes explicit that the conditional $P(D \mid G)$ is not required
to evaluate the Chow--Liu objective. The conditional may exist as a
regular conditional distribution, but it is not a standard local
component of the conditional Gaussian factorisation. Under an
additional sub-Gaussian deviation assumption on the pairwise mutual
information estimates with a common parameter $\sigma > 0$ (stated
formally as $P(|\hat{I}_n - I| > t) \le 2 \exp(-n t^2 / (2 \sigma^2))$
in \cref{cor:tree}), two finite-sample bounds follow. The first
is a tree-symmetric-difference bound (\cref{cor:tree}); the second
is a tighter per-swap bound (\cref{cor:swap}) that uses a per-edge
cycle-completion argument for Kruskal's algorithm and yields sample
complexity $O(\sigma^2 \log p / \gamma_{\mathrm{swap}}^2)$, where
$\gamma_{\mathrm{swap}}$ is the smallest population
mutual-information gap across all single-edge swaps in $T^*$
(formal definition at \eqref{eq:gamma_swap} and positivity under
(A2) in \cref{cor:swap}). Under (A2) the
per-swap gap coincides with the tree-symmetric-difference gap at
the population level (Remark after \cref{cor:swap}).

We do not establish the sub-Gaussian assumption itself.
Empirical variance-scale diagnostics for the continuous--continuous
KSG estimator and the mixed continuous--discrete $k$-NN MI
estimator \citep{gao2017estimating} are reported in
\cref{app:sigma}.

Before stating the forest analogue, we make three additional
assumptions explicit. Let $F^*$ denote the population thresholded
maximum spanning forest, that is, the forest obtained from Kruskal's
algorithm on the population MI matrix after discarding all edges
with $I(X_i; X_j) = 0$. The sample-level estimator uses a
shrinking threshold $\tau_n > 0$ in place of the population
zero-threshold rule, and the two thresholds agree only in the
limit $\tau_n \to 0$. We assume:
\begin{itemize}
\item[(A2$'$)] The population thresholded maximum spanning forest
$F^*$ is unique, equivalently the maximum spanning tree on each
connected component of the support graph
$\{(i, j) : I(X_i; X_j) > 0\}$ is unique.
\item[(A4)] Every edge in $F^*$ has mutual information at least
$\delta > 0$.
\item[(A5)] For every pair $(i, j)$ whose population mutual
information is zero, the estimator satisfies
$P\bigl(|\hat{I}_n(X_i; X_j)| > \tau_n\bigr) \to 0$
as $n \to \infty$.
\end{itemize}

\begin{proposition}[Forest consistency under shrinking-threshold rule]
\label{prop:forest}
Under (A1), (A3), (A2$'$), (A4), (A5), with $\tau_n \to 0$ and
$\tau_n \sqrt{n} \to \infty$, the estimated forest $\hat{F}_n$
obtained from Kruskal's algorithm with threshold $\tau_n$ satisfies
$P(\hat{F}_n = F^*) \to 1$.
\end{proposition}

The main forest experiments use $\tau_n = 0$ to match Edwards'
procedure, which uses no threshold. The shrinking-threshold
condition of \cref{prop:forest} ($\tau_n \to 0$, $\tau_n \sqrt{n}
\to \infty$) is therefore not formally satisfied in those
experiments. \Cref{subsec:shrinking} reports a separate experiment
that does satisfy the condition.

\subsection{The unrestricted Chow--Liu algorithm (Algorithm~1)}
\label{subsec:algorithm}

\begin{algBox}{Unrestricted mixed-type Chow--Liu algorithm}
\label{alg:main}
\textbf{Input:} Variables $X_1, \ldots, X_p$ with types
$t_i \in \{c, d\}$ ($c$ = continuous, $d$ = discrete); data
$\{x^{(1)}, \ldots, x^{(n)}\}$; threshold $\tau_n \ge 0$ (default
$\tau_n = 0$); neighbour count $k \in \mathbb{N}$ for the $k$-NN
estimators (default $k = 5$).
\begin{enumerate}
\item For each pair $(i, j)$ with $i < j$, compute the pairwise
mutual-information estimate according to the pair's type:
\begin{itemize}
\item if $t_i = c$ and $t_j = c$:
$\hat{w}_{ij} \gets \hat{I}^{(k)}_{\mathrm{KSG}}(X_i; X_j)$
\citep{kraskov2004estimating};
\item if $t_i = d$ and $t_j = d$:
$\hat{w}_{ij} \gets \hat{I}_{\mathrm{plug\text{-}in}}(X_i; X_j)$;
\item otherwise (one continuous and one discrete):
$\hat{w}_{ij} \gets \hat{I}^{(k),\,\mathrm{mix}}_{k\text{-NN}}(X_i; X_j)$
\citep{gao2017estimating}.
\end{itemize}
\item Exclude every edge with $\hat{w}_{ij} \le \tau_n$.
\item $\hat{F}_n \gets \mathrm{MaxSpanningForest}(\{\hat{w}_{ij}\})$,
computed by Kruskal's algorithm with no type restriction.
\end{enumerate}
\textbf{Return:} $\hat{F}_n$ (a spanning tree on $p$ vertices when
$\tau_n = 0$ and no edge is excluded, a strict forest otherwise).
\end{algBox}

\paragraph{Algorithm 1b (Gaussian variant).} For computational
efficiency on large graphs with many continuous--continuous pairs,
the closed-form Gaussian MI $\hat{I}_{cc}(X; Y) = -\tfrac{1}{2}\log(1
- \hat{\rho}^2)$ can replace the KSG estimator. We refer to this
variant as Algorithm~1b. It is faster but is consistent only under
joint Gaussianity. We use Algorithm~1b only in the breast-cancer
illustration (\cref{subsec:breastcancer}).

\paragraph{Algorithm 1a (Edwards-restricted baseline).} Throughout
the experiments we refer to Algorithm 1a as the Edwards-restricted
version of Algorithm 1, in which Kruskal's loop rejects any
candidate edge that would place a continuous node $X_c$ on the
unique path between two existing discrete components. Algorithm 1a
uses the same MI estimators and the same Kruskal ordering as
Algorithm 1. Only the type-check that decides whether to accept a
candidate edge differs. The empirical comparisons in
\cref{sec:experiments} use Algorithm 1a as the restricted baseline.

\paragraph{What Algorithm~1 provides relative to the
MDL-discretisation procedure of \citet{suzuki2017chow}.} Both
procedures lift Edwards' restriction and admit $D$-$G$-$D$ paths.
Algorithm~1 differs from the MDL-discretisation procedure in four
respects.
(i) It operates on the continuous data directly, with no information
loss due to discretisation (by the data-processing inequality,
$I(\tilde{G}; D) \le I(G; D)$ for any binning $\tilde{G}$ of $G$, with
equality only as the bin count tends to infinity).
(ii) It has a single hyperparameter ($k$, default 5), rather than a
bin scheme, a bin-count grid, and an MDL penalty constant.
(iii) \cref{prop:main} applies to any consistent pairwise MI
estimator, so Algorithm~1 is one instance among many, whereas the
MDL-discretisation procedure is tied to its specific discretisation
and penalty scheme.
(iv) \cref{prop:main} gives a direct consistency statement
for the Chow--Liu objective on the original continuous distribution,
whereas the MDL framework gives a description-length justification on
the discretised surrogate. The empirical comparison in
\cref{subsec:suzuki} shows that these methodological
differences translate into a sample-size-dependent practical
recommendation.

\subsection{Requirements on the mutual-information estimator}
\label{subsec:requirements}

\cref{prop:main} requires only per-pair consistency (A3).
The $k$-NN estimator
\citep{kraskov2004estimating, gao2017estimating} satisfies (A3) under
standard regularity conditions. The empirical plug-in for discrete
pairs is consistent. The Gaussian formula is consistent only under
joint Gaussianity. For variational mutual information estimators
(MINE, InfoNCE), \citet{mcallester2020formal} showed that achieving
$o(1)$ bias requires sample complexity exponential in the true mutual
information in the distribution-free setting. We therefore exclude
these estimators from the default procedure.

%%%%%%%%%%%%%%%%%%%%%%%%%%%%%%%%%%%%%%%%%%%%%%%%%%%%%%%%%%%%%%%%%%%%%%%%
\section{Experiments}
\label{sec:experiments}

All experiments use fixed random seeds and run on a single CPU core.
A single self-contained reproduction script
(\texttt{260516\_chow\_liu.py}) regenerates every numerical result and
every figure in the manuscript with the same publication-grade
settings used here. The script requires only NumPy, SciPy,
scikit-learn, pandas, matplotlib, and networkx. The full
reproduction script completes in approximately four minutes on a
single CPU core (laptop-grade hardware), and we report this figure
only as an order-of-magnitude indication.

We report two standard graph-recovery metrics throughout. The $F_1$
score is the harmonic mean of edge precision (the fraction of
recovered edges that are correct) and edge recall (the fraction of
true edges that are recovered). The value $F_1 = 1$ corresponds to
exact recovery. The structural Hamming distance (SHD) is the number
of edge insertions and deletions needed to turn the recovered tree
into the ground-truth tree; a lower SHD indicates better recovery.

\subsection{Accuracy of the mutual-information estimators}
\label{subsec:miacc}

We compare three estimators on four conditional distributions of
$G \mid D = d$, with $n = 2{,}000$ samples and 15 trials per scenario.
The true MI is computed by numerical integration. The conditional
families are Gaussian (location-shifted), bimodal (mixture of two
Gaussians), skewed (location-shifted $\chi^2_3$), and heavy-tailed
(location-shifted Student-$t_3$).

\begin{table}[ht]
\centering
\caption{MI estimation root mean squared error (RMSE), reported in
units of $10^{-2}$, on four conditional distributions ($n = 2{,}000$,
15 trials).}
\label{tab:mi_accuracy}
\begin{tabular}{lcccc}
\toprule
Scenario & True MI & ANOVA & Discret. & $k$-NN \\
\midrule
Gaussian   & 0.2538 & 1.13  & 3.38 & 1.55 \\
Bimodal    & 0.6053 & 36.80 & 8.44 & 2.15 \\
Skewed     & 0.1612 & 10.80 & 5.30 & 1.79 \\
Heavy-tail & 0.1738 & 8.09  & 3.01 & 1.93 \\
\bottomrule
\end{tabular}
\end{table}

ANOVA is optimal on Gaussian data but heavily biased on Bimodal
(RMSE about $17\times$ larger than that of $k$-NN) and on Skewed
(RMSE about $6\times$ larger). The $k$-NN estimator is competitive across all
four conditional families. We therefore use the $k$-NN estimator
(KSG for continuous--continuous pairs and the mixed-pair variant of
\citet{gao2017estimating} for continuous--discrete pairs) in the
empirical comparisons that follow, because the structure-learning
experiments cover heterogeneous-conditional regimes in which
ANOVA's linearity assumption does not hold.
\Cref{fig:accuracy} visualises the estimator distributions across
the four conditional families, and \cref{fig:convergence} shows
convergence of the three estimators with increasing sample size on
Gaussian data, the one benchmark on which ANOVA's assumption holds
exactly.

\begin{figure}[ht]
\centering
\includegraphics[width=0.9\textwidth]{Fig03_accuracy.png}
\caption{Mutual information estimation accuracy on four conditional
distributions of $G \mid D = d$ (Gaussian, bimodal, skewed,
heavy-tailed), each at $n = 2{,}000$ with 15 trials. The dashed line
is the true MI computed by numerical integration. The $k$-NN
estimator is competitive across all four families; ANOVA is severely
biased on the non-Gaussian families.}
\label{fig:accuracy}
\end{figure}

\begin{figure}[ht]
\centering
\includegraphics[width=0.8\textwidth]{Fig04_convergence.png}
\caption{Convergence of the three mutual-information estimators with
increasing sample size $n$ on the Gaussian benchmark. The dashed
line is the true MI. ANOVA has the smallest bias at every sample
size, because its linearity assumption holds exactly on this
benchmark. The $k$-nearest-neighbour estimator approaches the
truth from above with slightly larger small-sample bias than
ANOVA. The discretisation estimator has the largest small-sample
downward bias and the slowest approach to the true value, but
also approaches it as $n$ grows. All three estimators are
consistent on the Gaussian benchmark.}
\label{fig:convergence}
\end{figure}

\subsection{\texorpdfstring{$D$-$G$-$D$}{D-G-D} recovery (sanity check)}
\label{subsec:dgd_sanity}

We test the simplest possible $D$-$G$-$D$ configuration:
$D_1 \sim \mathrm{Uniform}\{0, 1, 2\}$, $G \sim \mathcal{N}(0, 1) +
1.5\, D_1$, and $D_2$ is a binary indicator of whether $G$ exceeds
its median. The true Markov structure is $D_1$-$G$-$D_2$, a path
that Edwards' restriction forbids by construction. In 50 trials at
$n = 1{,}500$ the unrestricted procedure recovers the path in
$100\%$ of trials, and the restricted procedure in $0\%$. The
difference is representational rather than statistical: even with
arbitrarily many samples, the restricted procedure cannot return
this structure. \Cref{fig:dgd} contrasts the structures the two
procedures can return.

\begin{figure}[ht]
\centering
\includegraphics[width=0.9\textwidth]{Fig05_dgd_chain.png}
\caption{Structural comparison on a $D_1$-$G$-$D_2$ chain. Top:
Algorithm~1 returns the true Markov path through the continuous
$G$. Bottom: Algorithm~1a, the Edwards-restricted baseline, cannot
represent the path through $G$ and instead returns the spurious
$D_1$-$D_2$ edge that skips $G$. This is a representational
constraint, not a sample-size effect. Recovery counts: Algorithm~1
$50 / 50 = 100\%$; Algorithm~1a $0 / 50 = 0\%$, by construction.}
\label{fig:dgd}
\end{figure}

\subsection{Synthetic expression quantitative trait locus (eQTL)
network}
\label{subsec:eqtl}

We consider an eight-variable expression quantitative trait locus
(eQTL) network with 3 SNPs, 4 gene expressions, and 1 disease
indicator, including the $D$-$G$-$D$ Markov path
$S_1 \to G_1 \to \text{Disease}$. We use $n = 1{,}500$ and 30 trials.

\begin{table}[ht]
\centering
\caption{Synthetic eQTL recovery, with structural Hamming distance
(SHD) defined as the number of edge insertions and deletions required
to transform the recovered tree into the ground-truth tree
($n = 1{,}500$, 30 trials).}
\label{tab:eqtl}
\begin{tabular}{lcccc}
\toprule
Method & Precision & Recall & $F_1$ & SHD \\
\midrule
Alg.~1 (proposed) + $k$-NN (KSG)        & 0.714 & 1.000 & 0.833 & 2.00 \\
Alg.~1a (Edwards-restr.) + $k$-NN (KSG) & 0.571 & 0.800 & 0.667 & 4.00 \\
\bottomrule
\end{tabular}
\end{table}

Every trial returns the same edge set; the variance across trials is
zero, so the recovery is deterministic at this signal level. The
network and the procedure outputs are shown in
\cref{fig:eqtl}.

\begin{figure}[ht]
\centering
\includegraphics[width=0.95\textwidth]{Fig06_eqtl.png}
\caption{Synthetic eQTL network. Discrete nodes (blue): SNPs
$S_1, S_2, S_3$; continuous nodes (green): gene expressions
$G_1, \ldots, G_4$; binary outcome (red): Disease. The red edges form
the $D$-$G$-$D$ Markov path $S_1 \to G_1 \to \text{Disease}$ that
Edwards' restriction forbids by construction. $F_1$ recovery and SHD
for each procedure are reported in \cref{tab:eqtl}.}
\label{fig:eqtl}
\end{figure}

\subsection{Random mixed graphs}
\label{subsec:randomgraphs}

We use a Pr\"ufer-sequence random-mixed-tree generator with rejection
sampling enforcing at least two $D$-$G$-$D$ paths in every trial. Data
are generated in breadth-first search (BFS) order from a randomly
chosen root, using the parent-to-child rules described in
\cref{app:datagen}. We use Algorithm~1 ($k$-NN throughout) on every
benchmark. Each table below draws an independent set of trials from
this generator with its own random seed, so the same nominal
configuration (for example, $p = 15$, $n = 1{,}500$, effect $0.6$)
gives slightly different mean $F_1$ across tables; the variation
across tables is therefore a rough indicator of generator-level
variability.

\paragraph{\texorpdfstring{(a) Scaling with $p$.}{(a) Scaling with p.}} We use
$n_c = \lfloor 8p/15 \rfloor$ continuous nodes, $n = 1{,}500$, 20
trials per row, and effect size 0.6.

\begin{table}[ht]
\centering
\caption{Scaling with $p$ on random mixed graphs ($n = 1{,}500$, 20
trials, effect size $= 0.6$).}
\label{tab:scaling}
\small
\begin{tabular}{cccccc}
\toprule
$p$ & $F_1$ Alg.~1 (prop.) & $F_1$ Alg.~1a (Edw.-restr.) & SHD diff. & paired $t$ & $P$ \\
\midrule
5   & $0.600 \pm 0.166$ & $0.438 \pm 0.108$ & $-1.3$ & $+4.95$ & $0.0001$ \\
8   & $0.679 \pm 0.201$ & $0.543 \pm 0.147$ & $-1.9$ & $+5.15$ & $0.0001$ \\
12  & $0.645 \pm 0.118$ & $0.573 \pm 0.108$ & $-1.6$ & $+3.76$ & $0.0013$ \\
15  & $0.739 \pm 0.149$ & $0.639 \pm 0.136$ & $-2.8$ & $+4.17$ & $0.0005$ \\
20  & $0.684 \pm 0.131$ & $0.637 \pm 0.110$ & $-1.8$ & $+2.39$ & $0.0273$ \\
30  & $0.728 \pm 0.079$ & $0.669 \pm 0.078$ & $-3.4$ & $+4.01$ & $0.0007$ \\
\bottomrule
\end{tabular}
\end{table}

\paragraph{(b) Effect-size robustness} ($p = 15$, $n = 1{,}500$, 20
trials).

\begin{table}[ht]
\centering
\caption{Effect-size sweep at $p = 15$.}
\label{tab:effect}
\begin{tabular}{ccccc}
\toprule
effect & $F_1$ Alg.~1 (prop.) & $F_1$ Alg.~1a (Edw.-restr.) & paired $t$ & $P$ \\
\midrule
0.4 & $0.575 \pm 0.127$ & $0.536 \pm 0.133$ & $+2.60$ & $0.0174$ \\
0.6 & $0.714 \pm 0.093$ & $0.650 \pm 0.110$ & $+4.16$ & $0.0005$ \\
0.8 & $0.746 \pm 0.129$ & $0.646 \pm 0.129$ & $+5.72$ & $<0.0001$ \\
1.0 & $0.768 \pm 0.115$ & $0.675 \pm 0.080$ & $+4.95$ & $0.0001$ \\
1.5 & $0.818 \pm 0.114$ & $0.689 \pm 0.094$ & $+9.00$ & $<0.0001$ \\
2.0 & $0.775 \pm 0.116$ & $0.661 \pm 0.095$ & $+4.77$ & $0.0001$ \\
\bottomrule
\end{tabular}
\end{table}

The paired advantage is significant across all effect sizes from
0.4 to 2.0. For interpretation, the true population mutual
information of a Gaussian-shift edge at each effect size is reported
in \cref{tab:effect_mi}. The regime spans true MI from
$\approx 0.05$ nats (effect $0.4$) to $\approx 0.62$ nats (effect
$2.0$).

\begin{table}[ht]
\centering
\caption{Population mutual information of a Gaussian-shift edge with
$K = 3$ discrete categories, unit conditional variance, at each
effect size used in \cref{tab:effect}. Computed by numerical
integration.}
\label{tab:effect_mi}
\begin{tabular}{cc}
\toprule
effect & true $I(G; D)$ (nats) \\
\midrule
0.4 & 0.0507 \\
0.6 & 0.1075 \\
0.8 & 0.1772 \\
1.0 & 0.2538 \\
1.5 & 0.4472 \\
2.0 & 0.6191 \\
\bottomrule
\end{tabular}
\end{table}

\paragraph{(c) Heterogeneous-distribution data} ($p = 15$,
$n = 1{,}500$, 20 trials, effect size $= 0.6$): each $c \mid d$ edge
independently uses one of four conditional families, chosen uniformly
at random. The families are Gaussian
$\mathcal{N}(\text{effect}\cdot d, 1)$, bimodal, skewed
($\chi^2_3$-shift), and heavy-tailed ($t_3$-shift).

\begin{table}[ht]
\centering
\caption{Heterogeneous-distribution experiment. The
MDL-discretisation MI estimator \citep{suzuki2017chow} is added as a
comparator; paired $t$-tests are against the proposed pipeline with
the KSG estimator.}
\label{tab:hetero}
\begin{tabular}{lcc}
\toprule
Method & $F_1$ & paired $t$ ($P$) vs Alg.~1 + KSG \\
\midrule
Alg.~1 (proposed) + $k$-NN (KSG)        & $0.707 \pm 0.098$ & reference \\
Alg.~1 (proposed) + MDL discr.\ MI      & $0.704 \pm 0.116$ & $+0.15$ ($P = 0.8833$) \\
Alg.~1a (Edwards-restr.) + $k$-NN (KSG) & $0.654 \pm 0.099$ & $+3.29$ ($P = 0.0039$) \\
\bottomrule
\end{tabular}
\end{table}

\Cref{tab:hetero} reports the result. Under heterogeneous
conditional families, Algorithm~1 paired with the KSG estimator and
Algorithm~1 paired with MDL discretisation yield indistinguishable
mean $F_1$ ($P = 0.88$). The proposed pipeline
significantly outperforms the Edwards-restricted baseline
($P = 0.0039$). The advantage of removing the type-pair restriction
therefore holds under distributional heterogeneity; the choice
between the two MI estimators shows no significant difference on
this benchmark at $n = 1{,}500$.

\subsection{Sample-size sweep on the
\texorpdfstring{$D$-$G$-$D$}{D-G-D} chain}
\label{subsec:neartie}

We revisit the three-node $D_1$-$G$-$D_2$ chain of
\cref{subsec:dgd_sanity} ($D_1 \sim \mathrm{Uniform}\{0,1,2\}$, $G$
Gaussian shifted by $D_1$, $D_2$ a thresholded function of $G$) and
vary the sample size. Chain transitivity gives the direct $D_1$-$D_2$ pair a non-zero
population mutual information that lies below the two true path
edges. A noisy MI estimate could therefore swap the true edge
$G$-$D_2$ for the spurious $D_1$-$D_2$ edge during the Kruskal
step, making this configuration a finite-sample stress test for
the unrestricted procedure. It
recovers the $D$-$G$-$D$ chain in 100\% of trials across
$n \in \{500, 1000, 2000, 5000\}$,
while the restricted procedure recovers 0\% by construction.

\subsection{Shrinking-threshold variant of Algorithm~1}
\label{subsec:shrinking}

We run Algorithm~1 with $\tau_n = c \log(n)/\sqrt{n}$ for
$c \in \{0, 0.1, 0.5, 1.0\}$, on $p = 15$ random mixed graphs from the
generator used in \cref{subsec:randomgraphs}, with $n =
1{,}500$, 20 trials, and effect size 0.6. The setting $c = 0$
corresponds to the zero-threshold rule used elsewhere, while
$c \in \{0.1, 0.5, 1.0\}$ satisfies the conditions of
\cref{prop:forest}.

\begin{table}[ht]
\centering
\caption{Shrinking-threshold variant.}
\label{tab:shrinking}
\begin{tabular}{cccccc}
\toprule
$c$ & $\tau_n$ at $n=1500$ & $F_1$ & SHD & Precision & Recall \\
\midrule
0.0 & 0.000 & $0.679 \pm 0.127$ & $9.0 \pm 3.6$ & 0.679 & 0.679 \\
0.1 & 0.019 & $0.679 \pm 0.127$ & $9.0 \pm 3.6$ & 0.679 & 0.679 \\
0.5 & 0.094 & $0.562 \pm 0.136$ & $10.4 \pm 3.0$ & 0.677 & 0.482 \\
1.0 & 0.189 & $0.466 \pm 0.086$ & $11.0 \pm 1.8$ & 0.731 & 0.343 \\
\bottomrule
\end{tabular}
\end{table}

At $n = 1{,}500$, the threshold $c = 0.1$ excludes no true edges. In
practice, at $n = 1{,}500$ the choice $c = 0.1$ yields $F_1$ and
SHD identical to the zero-threshold rule (\cref{tab:shrinking})
while satisfying the consistency conditions of
\cref{prop:forest}.

\subsection{Sample-size sweep against the MDL-discretisation estimator}
\label{subsec:suzuki}

The MDL-discretisation MI estimator of \citet{suzuki2017chow}
discretises each continuous variable and uses the MDL-penalised
empirical MI
$\hat{I}_{\mathrm{MDL}}(\tilde{G}; D) = \hat{I}(\tilde{G}; D) -
(M-1)(C-1)\log(n)/(2n)$. We implement this estimator by sweeping bin
counts $M \in \{2, 4, 8, 16, 32\}$ and taking the maximum penalised
MI, and feed its output into the same unrestricted Kruskal step as
Algorithm~1.

\paragraph{Synthetic eQTL benchmark} ($n = 1{,}500$, 30 trials).
Algorithm~1 paired with the MDL-discretisation MI estimator and
Algorithm~1 paired with the KSG estimator produce identical
results. Both yield $F_1 = 0.833$, SHD $= 2.0$, and both recover the
$S_1 \to G_1 \to \text{Disease}$ chain. The Edwards-restricted
baseline yields $F_1 = 0.667$ and SHD $= 4.0$.

\paragraph{Random mixed graphs benchmark, sample-size sweep}
($p = 15$, 20 trials per row, effect size $= 0.6$).

\begin{table}[ht]
\centering
\caption{Sample-size sweep on $p = 15$ random graphs.}
\label{tab:suzuki_n}
\small
\begin{tabular}{cccc}
\toprule
$n$ & $F_1$ Alg.~1 + KSG & $F_1$ Alg.~1 + MDL discr. & $F_1$ Alg.~1a (Edw.-restr.) + KSG \\
\midrule
75    & $0.529 \pm 0.086$         & $\mathbf{0.575 \pm 0.107}$ & $0.521 \pm 0.085$ \\
100   & $0.568 \pm 0.136$         & $0.542 \pm 0.116$          & $0.521 \pm 0.132$ \\
200   & $0.625 \pm 0.087$         & $0.660 \pm 0.091$          & $0.586 \pm 0.080$ \\
500   & $0.589 \pm 0.133$         & $0.619 \pm 0.106$          & $0.543 \pm 0.120$ \\
1000  & $0.707 \pm 0.093$         & $0.664 \pm 0.099$          & $0.632 \pm 0.079$ \\
1500  & $\mathbf{0.736 \pm 0.122}$ & $0.686 \pm 0.104$          & $0.657 \pm 0.086$ \\
\bottomrule
\end{tabular}
\end{table}

We report paired $t$-tests on $F_1$ over the same simulated graphs.
The labels indicate which procedure has the larger mean $F_1$ when the
test is significant at $\alpha = 0.05$; otherwise the entry is marked
``n.s.''\ (not significant).

\begin{table}[ht]
\centering
\caption{Paired comparisons for the sample-size sweep.}
\label{tab:suzuki_ttest}
\small
\begin{tabular}{cc}
\toprule
$n$ & KSG vs MDL discr.\ MI: paired $t$-test on $F_1$ \\
\midrule
75    & MDL discr.\ preferred, $P = 0.0395$ \\
100   & n.s., $P = 0.2308$ \\
200   & n.s., $P = 0.1885$ \\
500   & n.s., $P = 0.2697$ \\
1000  & n.s., $P = 0.0967$ (KSG slightly higher mean) \\
1500  & KSG preferred, $P = 0.0265$ \\
\bottomrule
\end{tabular}
\end{table}

A qualitative regime crossover is visible between the two
extremes. Denote the two pairings as Alg.~1\,+\,MDL and
Alg.~1\,+\,KSG. At $n = 75$, Alg.~1\,+\,MDL yields significantly
higher $F_1$ ($P = 0.04$), because the MDL penalty adaptively
favours coarser bins when samples are sparse, which acts as a
form of model selection. At $n \in \{100, 200, 500\}$ the two
estimators are statistically indistinguishable on $F_1$
($P \in \{0.23, 0.19, 0.27\}$ respectively). At $n = 1{,}000$,
Alg.~1\,+\,KSG has slightly higher mean $F_1$, but the paired
comparison is not significant at $\alpha = 0.05$ ($P = 0.097$).
At $n = 1{,}500$, Alg.~1\,+\,KSG is significantly better
($P = 0.027$). The discretisation step in the MDL estimator
introduces a finite-sample bias that grows relative to its
variance benefit as $n$ increases. (The
data-processing inequality says only that
$I(\tilde G; D) \le I(G; D)$ at the population level; at finite
sample sizes, discretisation is a bias--variance trade-off, and the
MDL penalty adaptively selects coarser bins to favour variance
reduction under data sparsity.) At both ends of the sample-size
range the unrestricted Algorithm~1 yields higher mean $F_1$ than the
Edwards-restricted baseline.

\paragraph{\texorpdfstring{Adaptive-$k$ variant (exploratory).}{Adaptive-k variant (exploratory).}}
An exploratory adaptive-$k$ check is reported in
\cref{app:adaptive}; it narrows but does not eliminate the
small-sample advantage of the MDL-discretisation estimator, and
is not used in support of the main $n = 1{,}500$ comparison.

\subsection{Mutual-information estimator ablation}
\label{subsec:ablation}

To check whether the recovery improvement reported in
\cref{subsec:eqtl} depends on the specific MI estimator, we
compare three estimators (ANOVA, fixed-grid discretisation, and
$k$-NN (KSG)) on the same eQTL benchmark at $n = 1{,}500$ over
20 trials. All three are used as inputs to Algorithm~1.

\begin{table}[ht]
\centering
\caption{Estimator ablation on the eQTL benchmark.}
\label{tab:ablation}
\begin{tabular}{lcc}
\toprule
Procedure & $F_1$ & SHD \\
\midrule
Alg.~1 (proposed) + ANOVA MI               & $0.833 \pm 0.000$ & $2.00 \pm 0.00$ \\
Alg.~1 (proposed) + fixed-grid discr.\ MI  & $0.833 \pm 0.000$ & $2.00 \pm 0.00$ \\
Alg.~1 (proposed) + $k$-NN (KSG) MI        & $0.833 \pm 0.000$ & $2.00 \pm 0.00$ \\
Alg.~1a (Edwards-restr.) + $k$-NN (KSG) MI & $0.667 \pm 0.000$ & $4.00 \pm 0.00$ \\
\bottomrule
\end{tabular}
\end{table}

On this benchmark, all three MI estimators produce identical
recovery. The structural advantage of removing the restriction
(0.833 versus 0.667 in $F_1$) is independent of the choice. We
recommend $k$-NN as the default.

\subsection{Fully discrete sanity check: Asia network}
\label{subsec:asia}

We include the Asia Bayesian network \citep{lauritzen1988local} as a
sanity check on a canonical fully discrete benchmark. Asia has 8
binary variables and 8 directed acyclic graph (DAG) skeleton edges. We use the 5000-sample dataset from the
\texttt{gRbase} R package \citep{hojsgaard2012graphical}. For
$n \in \{500, 1{,}000, 2{,}500\}$, we draw 15 subsamples without
replacement from the full 5000-sample dataset; the $n = 5{,}000$
row uses the full dataset directly. Asia is
fully discrete, so Edwards' type-pair restriction is inactive
and all three procedures (Alg.~1 + KSG, Alg.~1a + KSG, Alg.~1 + MDL
discretisation) produce identical trees by construction. The
benchmark therefore serves to confirm that the proposed pipeline
recovers a known canonical structure on real data, rather than to
provide a comparison among methods.
The Asia skeleton has 8 edges, but any spanning tree on 8 nodes
contains only 7 edges, so a Chow--Liu tree cannot recover the full
skeleton. The maximum $F_1$ attainable by a tree approximation of
an 8-node, 8-edge graph occurs when all 7 tree edges are true
positives: precision $= 7/7 = 1$, recall $= 7/8$, and
$F_1 = \tfrac{2 \cdot 1 \cdot (7/8)}{1 + 7/8} = 14/15 \approx 0.933$.
\Cref{tab:asia} reports exactly this value at $n = 5{,}000$, so the
procedure attains the tree-approximation ceiling.

\begin{table}[ht]
\centering
\caption{Asia network skeleton recovery (the three procedures coincide
on this benchmark).}
\label{tab:asia}
\begin{tabular}{cc}
\toprule
$n$ & $F_1$ \\
\midrule
500   & $0.803 \pm 0.091$ \\
1000  & $0.800 \pm 0.069$ \\
2500  & $0.827 \pm 0.053$ \\
5000  & $0.933$ (full sample, single run) \\
\bottomrule
\end{tabular}
\end{table}

At $n = 5{,}000$, the procedure recovers 7 of the 8 DAG skeleton edges
in every subsample. This matches the maximum number of skeleton edges
that any spanning tree on 8 nodes can contain (a tree has $p - 1 = 7$
edges); the missing edge is excluded by the tree edge-count
constraint, not by estimation failure.

\subsection{Breast cancer gene expression (illustration only)}
\label{subsec:breastcancer}

We apply Algorithm~1b to $n = 250$ samples and 1000 genes
\citep{west2001predicting}. The top gene identified by both the
$t$-test $p$-value and the $k$-NN mutual information against the
class label is \texttt{A.202870\_s\_at}.

\begin{table}[ht]
\centering
\caption{Top five genes ranked by $k$-NN mutual information against
the class label.}
\label{tab:bc}
\begin{tabular}{clr}
\toprule
Rank & Probe & $k$-NN MI \\
\midrule
1 & A.202870\_s\_at & 0.236 \\
2 & A.212949\_at    & 0.214 \\
3 & A.209408\_at    & 0.208 \\
4 & B.224753\_at    & 0.190 \\
5 & B.228559\_at    & 0.189 \\
\bottomrule
\end{tabular}
\end{table}

No ground-truth graph is available for this dataset. The section is
included for illustration only. Algorithm~1b uses the Gaussian MI
formula for gene--gene pairs. \Cref{fig:bc_forest} shows the
top-50-gene forest learned by the proposed procedure on this dataset.

\begin{figure}[ht]
\centering
\includegraphics[width=0.95\textwidth]{Fig08_forest50.png}
\caption{Chow--Liu forest learned by Algorithm~1b on the
breast-cancer expression dataset (top 50 differentially expressed
genes plus the binary class label, $n = 250$). Gene--class edges use
the mixed-pair $k$-NN mutual information estimator; gene--gene edges
use the Gaussian formula. The class node (CLASS, red) is connected
to gene \texttt{A.202870\_s\_at}, which is also the top gene by
$t$-test $p$-value. The section is presented as an illustration; no
ground-truth graph is available for this dataset.}
\label{fig:bc_forest}
\end{figure}

%%%%%%%%%%%%%%%%%%%%%%%%%%%%%%%%%%%%%%%%%%%%%%%%%%%%%%%%%%%%%%%%%%%%%%%%
\section{Discussion}
\label{sec:discussion}

The propositions of \cref{subsec:consistency} formalise the
observation that Edwards' restriction is unnecessary at the
structure-selection step. \Cref{cor:swap} is the main quantitative
contribution. It yields a sample complexity of
$O(\sigma^2 \log p / \gamma_{\mathrm{swap}}^2)$. Under (A2),
$\gamma_{\mathrm{swap}} = \gamma$ at the population level (because
the second-best spanning tree differs from $T^*$ by a single edge
swap), and \cref{cor:swap} is tighter than \cref{cor:tree} by a
factor of $(p-1)^2$ in the exponent.

Empirically, on the $n = 1{,}500$ graph-size sweep with random mixed
trees containing $D$-$G$-$D$ Markov paths, removing the restriction
yields higher $F_1$ at every tested graph size, with paired tests
remaining significant after Benjamini--Hochberg correction
(\cref{subsec:randomgraphs}). On the eQTL benchmark
(\cref{subsec:eqtl}) the advantage is $F_1 = 0.833$ versus
$0.667$; the recovery is deterministic at this signal level. Against the MDL-discretisation MI estimator of
\citet{suzuki2017chow} (\cref{subsec:suzuki}), the comparison is
sample-size dependent. At $n = 75$, pairing Algorithm~1 with MDL
discretisation yields significantly higher $F_1$ than pairing it
with KSG ($P = 0.04$). At $n \in \{100, 200, 500, 1{,}000\}$ the
two estimators are statistically indistinguishable
($P > 0.05$ in every case). At $n = 1{,}500$, pairing with KSG is
significantly better ($P = 0.027$). On the fully discrete Asia
Bayesian network (\cref{subsec:asia}) all three procedures (Algorithm
1 with KSG, Algorithm 1a with KSG, Algorithm 1 with MDL
discretisation) coincide because Edwards' restriction is inactive on
all-discrete data; at $n = 5{,}000$, the learned tree recovers $7$ of
the $8$ skeleton edges, the maximum possible number for any spanning
tree on $8$ nodes.

\paragraph{Multiple testing.} \cref{subsec:randomgraphs,subsec:suzuki}
report $20$ paired $t$-tests across the scaling, effect-size,
heterogeneous, and MI-estimator sample-size sweeps. Of these,
$15$ have raw $P < 0.05$ and $5$ have raw $P \ge 0.05$. After applying
the Benjamini--Hochberg procedure to control the false discovery
rate at $0.05$, we find that $14$ of the $15$ raw-significant
tests remain significant after correction (largest adjusted $P = 0.039$); the
remaining raw-significant comparison (MDL versus KSG at $n = 75$,
raw $P = 0.0395$) lies above the BH threshold (adjusted
$P = 0.053$). The five raw non-significant comparisons remain
non-significant under the correction. Multiple-testing adjustment therefore does not change the
qualitative claim.

\Cref{app:sigma} reports an empirical variance-scale
diagnostic for the KSG and mixed-pair $k$-NN estimators. The
diagnostic returns the quantity
$\sigma_{\mathrm{emp}} = \sqrt{n \cdot \mathrm{Var}(\hat{I}_n)}$,
which is the empirical analogue of the sub-Gaussian parameter
$\sigma$ from \cref{cor:tree}. The value $\sigma_{\mathrm{emp}} = 1.2$ is an upper envelope across
all Gaussian-pair correlations and Gaussian-shift effect sizes
tested in \cref{app:sigma}; the largest cell in the underlying tables is $1.10$ for the
continuous--continuous KSG estimator and $0.89$ for the
mixed-pair $k$-NN estimator, and we round upward to $1.2$
(approximately a $9\%$ margin above the maximum observed value). The diagnostic is
useful for interpreting the scale of fluctuations in the
experiments, but it does not establish the sub-Gaussian deviation
assumption used in \cref{cor:tree,cor:swap}. The numerical substitution $\sigma = 1.2$ into the
sample-complexity formula should therefore be read as
illustrative under the assumed tail bound, not as a proven
finite-sample guarantee. A formal derivation of $\sigma$ as a
function of explicit smoothness assumptions on the densities, in the
spirit of \citet{singh2014exponential}, \citet{berrett2019efficient},
and \citet{han2020optimal}, would convert the empirical diagnostic
into a theorem.

Our bounds depend on a mutual-information gap (the
tree-symmetric-difference gap $\gamma$ in \cref{cor:tree} and the
per-swap gap $\gamma_{\mathrm{swap}}$ in \cref{cor:swap}, which
coincide under (A2)), but we do not derive these gaps from more
interpretable signal properties such as effect sizes or
conditional variances. A future analysis that expressed
$\gamma_{\mathrm{swap}}$ in terms of such quantities would tie the
sample-complexity formula directly to features of the
data-generating model.

\paragraph{Practical reach of the per-swap bound on the random-tree
generator.} An empirical proxy for $\gamma_{\mathrm{swap}}$ computed on the
estimated population MI matrix of the random-mixed-tree generator
(\cref{app:datagen}) at $n = 8{,}000$ has median $\approx 0$ nats
(IQR $[0, 0.001]$ and maximum $0.0035$ nats across 15 trials). At
the median, (A2) effectively fails and the per-swap bound is
undefined. Even at the largest $\gamma_{\mathrm{swap}}$ observed
across the 15 trials, substituting $\gamma_{\mathrm{swap}} = 0.0035$
nats into the \cref{cor:swap} sample-complexity formula at $\sigma
= 1.2$, $p = 20$, $\delta = 0.1$ yields $n \ge 138 / (0.0035)^2
\approx 1.1 \times 10^{7}$. \Cref{cor:swap} is therefore not
numerically informative on the random-mixed-tree benchmark used in
\cref{subsec:randomgraphs,subsec:suzuki}; the content of the bound
is diluted by near-ties in the population MI matrix produced by
that generator.

\paragraph{\texorpdfstring{A bounded-$\gamma_{\mathrm{swap}}$ regime where \cref{cor:swap} is numerically informative.}{A bounded-gamma\_swap regime where Corollary 2 is numerically informative.}} As a
controlled demonstration, we consider a three-node $D_1$-$G$-$D_2$
Markov chain with strong effect: $D_1 \sim \mathrm{Uniform}\{0,1,2\}$,
$G = 2 D_1 + \varepsilon$ with $\varepsilon \sim \mathcal{N}(0, 1)$,
and $D_2$ obtained by binning $G$ at its population terciles. The
population MIs are obtained in closed form except for $H(G)$, which
is computed by one-dimensional quadrature against the Gaussian
mixture $(1/3)\sum_{d=0}^{2} \mathcal{N}(2d, 1)$. The conditional
entropy $H(G \mid D_1)$ equals $\tfrac{1}{2}\log(2\pi e)$, since each
conditional density is $\mathcal{N}(2d, 1)$. The joint $P(D_1, D_2)$
is a sum of three normal-CDF differences, giving $H(D_1, D_2)$ in
closed form. The mutual information $I(G; D_2) = \log 3$ because
$D_2$ is a deterministic function of $G$ at the population
terciles. This gives
$I(D_1; G) = 0.6191$ nats, $I(G; D_2) = 1.0986$ nats,
$I(D_1; D_2) = 0.5202$ nats, and a per-swap gap
$\gamma_{\mathrm{swap}} = 0.0990$ nats. Substituting into
\cref{cor:swap} at $\sigma_{\mathrm{emp}} = 1.2$, $p = 3$,
$\delta = 0.1$ gives a required sample size of $n \approx 7{,}400$;
an empirical evaluation over $100$ trials at $n = 1{,}500$ achieves
$100\%$ recovery on this regime. The bound is therefore conservative
but numerically finite when $\gamma_{\mathrm{swap}}$ is bounded
away from zero.

\paragraph{\texorpdfstring{Extension to $p = 15$.}{Extension to p = 15.}}
To verify that the same finite-sample story holds beyond the
three-node case, we form a $p = 15$ generator by stacking five
independent copies of the above $D_1$-$G$-$D_2$ chain. The
population MI matrix is block-diagonal, so the per-block
$\gamma_{\mathrm{swap}} = 0.0990$ nats from the three-node case
carries over. Substituting into \cref{cor:swap} at $\sigma =
1.2$, $p = 15$, $\delta = 0.1$ gives a required sample size of
$n \approx 13{,}100$. The truth is a $10$-edge forest, whereas the
maximum-spanning-tree step returns a $14$-edge spanning tree on the
$15$ nodes; we therefore report whether all ten true edges are
included in the returned spanning tree. Across $100$ trials at
$n = 1{,}500$, all ten true edges are present in the returned tree
in $100\%$ of trials, with edge-set $F_1 = 0.833$.\footnote{Because
across-block population MI is exactly zero in this construction,
assumption (A2) of unique-maximum MST does not hold globally. The
demonstration is therefore most cleanly viewed as an instance of
\cref{prop:forest}, the thresholded-forest version: under a
shrinking threshold $\tau_n$ that excludes near-zero MI edges, the
returned forest is exactly the ten-edge truth and edge-set $F_1$
tends to $1$. We report the spanning-tree statistic above to keep
the comparison directly aligned with the three-node demonstration in
the preceding paragraph.} The bound's prediction is satisfied well
before the formal required sample size, as expected of a
conservative concentration bound. The general message is that \cref{cor:swap} is informative
for problems whose population MI matrix has separated edge weights,
and that it becomes numerically vacuous when the population matrix
admits near-ties; the
random-mixed-tree generator used in \cref{subsec:randomgraphs} and
\cref{subsec:suzuki} produces the latter regime.

%%%%%%%%%%%%%%%%%%%%%%%%%%%%%%%%%%%%%%%%%%%%%%%%%%%%%%%%%%%%%%%%%%%%%%%%
\section{Conclusion}
\label{sec:conclusion}

We have revisited the mixed-type Chow--Liu algorithm and asked
whether Edwards' type-pair restriction, which forbids a
continuous variable from connecting two discrete variables along
a tree path, is needed at the structure-selection stage. We
showed that it is not. Tree selection depends only on pairwise
mutual-information values, each symmetric in its two arguments
and well-defined for every combination of variable types, so the
Chow--Liu objective does not require a directional local factor
such as $P(D \mid G)$. The restriction enters only at the
subsequent step that fits a parametric density on the chosen
tree.

\Cref{prop:main} shows that the unrestricted procedure recovers
the true Chow--Liu tree as the sample size grows, whenever the
pairwise mutual-information estimator is per-pair consistent.
\Cref{prop:forest} extends this conclusion to a thresholded
forest analogue under a shrinking-threshold rule.
\Cref{cor:tree,cor:swap} supply finite-sample tree-recovery
bounds under a per-pair sub-Gaussian deviation assumption with a
common parameter $\sigma$, and the per-swap bound in
\cref{cor:swap} sharpens the tree-symmetric-difference bound in
\cref{cor:tree} by a factor of $(p-1)^2$ in the exponent.
\Cref{app:sigma} reports a variance-scale diagnostic that
supports the sub-Gaussian substitution numerically without proving
the tail itself.

Empirically, the unrestricted procedure outperforms the
Edwards-restricted baseline on random mixed graphs containing
$D$-$G$-$D$ Markov paths that the restricted procedure cannot
represent. Of the $20$ paired $t$-tests reported in
\cref{subsec:randomgraphs,subsec:suzuki}, $15$ are significant at
$\alpha = 0.05$ and $14$ remain significant under
Benjamini--Hochberg correction at the same level. The synthetic eQTL
network corroborates the recovery improvement on a biologically
motivated structure, and the fully discrete Asia Bayesian
benchmark, where Edwards' restriction is inactive by
construction, serves as a separate sanity check that the
pipeline attains the tree-approximation ceiling on a canonical
discrete network. Substituting the MDL-discretisation estimator for the
$k$-NN estimator inside the same unrestricted pipeline produces a
sample-size-dependent comparison. MDL discretisation outperforms
$k$-NN at $n = 75$; the two estimators perform indistinguishably
at intermediate sample sizes ($n \in \{100, 200, 500, 1{,}000\}$);
and $k$-NN outperforms MDL discretisation at $n = 1{,}500$.

Two caveats apply. First, the returned graph is undirected, and
its edges represent statistical associations rather than causal
relations. Second, the result covers structure selection only.
If a later step fits a parametric density on the chosen tree,
Edwards' restriction or an alternative parametrisation of the
offending local conditional becomes relevant again at that step,
and we do not address it here.

%%%%%%%%%%%%%%%%%%%%%%%%%%%%%%%%%%%%%%%%%%%%%%%%%%%%%%%%%%%%%%%%%%%%%%%%
\section*{Funding}

No funding was received for this project.

%%%%%%%%%%%%%%%%%%%%%%%%%%%%%%%%%%%%%%%%%%%%%%%%%%%%%%%%%%%%%%%%%%%%%%%%
\section*{Conflict of interest}

The authors declare no conflict of interest.

%%%%%%%%%%%%%%%%%%%%%%%%%%%%%%%%%%%%%%%%%%%%%%%%%%%%%%%%%%%%%%%%%%%%%%%%
\section*{Code and data availability}

The full reproduction code, including the random-mixed-tree
generator, the experiment scripts, and the figure-generation
pipeline, is available at the repository
\begin{center}
\url{https://github.com/ash141886/Unrestricted-Chow-Liu-}.
\end{center}
The Asia network dataset is included with the \texttt{gRbase} R
package \citep{hojsgaard2012graphical}, and the breast-cancer
gene-expression dataset of \citet{west2001predicting} is publicly
available.

%%%%%%%%%%%%%%%%%%%%%%%%%%%%%%%%%%%%%%%%%%%%%%%%%%%%%%%%%%%%%%%%%%%%%%%%
\section*{Declaration on the use of generative AI}

During the preparation of this manuscript, the authors used a
large language model (ChatGPT) to assist with writing polish,
notation consistency checks, and LaTeX formatting. All
theoretical results, derivations, experimental designs, software
implementations, and final scientific claims were originated and
verified by the authors. The authors reviewed and edited the
AI-assisted content and take full responsibility for the content
of the publication.

%%%%%%%%%%%%%%%%%%%%%%%%%%%%%%%%%%%%%%%%%%%%%%%%%%%%%%%%%%%%%%%%%%%%%%%%
\section*{Appendix}

The appendix contains the proofs of \cref{prop:main,prop:forest},
the proofs of \cref{cor:tree,cor:swap}, the data-generation rules
used by the random-mixed-tree generator, the empirical
variance-scale diagnostic supporting the sub-Gaussian
substitutions used in \cref{cor:tree,cor:swap}, and an exploratory
adaptive-$k$ table.

\appendix

\section{Proofs and finite-sample bounds}
\label{app:proofs}

\subsection*{Proof of \cref{prop:main}}

Define the population objective
$W(T) = \sum_{(i,j) \in T} I(X_i; X_j)$ and its sample analogue
$\hat{W}_n(T) = \sum_{(i,j) \in T} \hat{I}_n(X_i; X_j)$. The Chow--Liu
estimator is $\hat{T}_n = \arg\max_{T \in \mathcal{T}} \hat{W}_n(T)$,
and the target is $T^* = \arg\max_{T \in \mathcal{T}} W(T)$.

\emph{Step 1 (uniform convergence over $\mathcal{T}$).} Every spanning
tree has exactly $p - 1$ edges, so
\begin{equation*}
\sup_{T \in \mathcal{T}} |\hat{W}_n(T) - W(T)| \le (p-1)
\max_{i < j} |\hat{I}_n(X_i; X_j) - I(X_i; X_j)|
\overset{p}{\to} 0,
\end{equation*}
where the convergence follows from (A3) applied to each of the
$\binom{p}{2}$ pairs.

\emph{Step 2 (identification).} By (A2), $T^*$ is unique, so there
exists a positive gap
$\gamma = W(T^*) - \max_{T \ne T^*} W(T) > 0$.

\emph{Step 3 (consistency).} For any $T \ne T^*$,
\begin{equation*}
\hat{W}_n(T^*) - \hat{W}_n(T) \;\ge\; \gamma - 2 \sup_{T'}
|\hat{W}_n(T') - W(T')|.
\end{equation*}
The right-hand side exceeds $\gamma/2 > 0$ with probability tending
to 1, so $P(\hat{T}_n = T^*) \to 1$. This is the standard M-estimation
argument \citep{vandervaart2000asymptotic}. \qed

\medskip
\noindent\emph{Remark (no local CG parametrisation $P(D \mid G)$ is
required).} For a mixed pair $(G, D)$, mutual information is
symmetric:
\begin{equation}
I(G; D) \;=\; H(G) - H(G \mid D) \;=\; H(D) - H(D \mid G).
\label{eq:identity}
\end{equation}
The Chow--Liu objective depends only on the scalar pairwise quantity
$I(G; D)$, not on a chosen directed local factorisation of the pair.
Evaluating or estimating an edge weight therefore does not require
specifying a standard conditional-Gaussian local factor $P(D \mid
G)$. Such a factor may be required at a later, separate step that
writes the joint distribution as an explicit parametric density on
the selected tree.

\subsection{Data-generation rules for the random-mixed-tree
generator}
\label{app:datagen}

The random-mixed-tree generator uses a Pr\"ufer sequence to sample a
labelled tree uniformly. Conditional on the tree, data are generated
in BFS order from a randomly chosen root, using the following
parent-to-child rules:

\begin{itemize}
\item Discrete parent $D$, continuous child $G$:
$G = \text{effect} \cdot D + \varepsilon$, with
$\varepsilon \sim \mathcal{N}(0, 1)$.
\item Continuous parent $G$, discrete child $D$:
$D = \mathbf{1}[G > q_{1/3}] + \mathbf{1}[G > q_{2/3}]$, where
$q_{1/3}$ and $q_{2/3}$ are the empirical tercile cutoffs of $G$.
\item Two discrete: noisy copy with a 30\% flip rate.
\item Two continuous:
$G_{\mathrm{child}} = \alpha G_{\mathrm{parent}} + 0.8 \varepsilon$,
with $\alpha = \text{effect}/(1 + \mathrm{sd}(G_{\mathrm{parent}}))$.
\end{itemize}

\subsection{Tree-symmetric-difference bound}
\label{app:cor1}

\begin{corollary}[Tree-symmetric-difference bound]
\label{cor:tree}
Assume (A1)--(A2) and that, for every pair $(i, j)$ with a common
parameter $\sigma > 0$, the pairwise estimator satisfies the
sub-Gaussian tail
\begin{equation}
P\bigl(|\hat{I}_n(X_i; X_j) - I(X_i; X_j)| > t\bigr) \;\le\; 2 \exp\!\left(-\frac{n t^2}{2 \sigma^2}\right), \qquad t > 0.
\label{eq:subgauss}
\end{equation}
This tail implies (A3) and so subsumes the consistency hypothesis of
\cref{prop:main}. With
$\gamma = W(T^*) - \max_{T \ne T^*} W(T) > 0$,
\begin{equation}
P(\hat{T}_n \ne T^*) \;\le\; p^2 \exp\!\left(-\frac{n \gamma^2}{8
\sigma^2 (p-1)^2}\right).
\end{equation}
\end{corollary}

\emph{Proof.} From the uniform-bound step of \cref{prop:main} (Step
1), $\hat{T}_n = T^*$ whenever
$\sup_{T \in \mathcal{T}} |\hat{W}_n(T) - W(T)| < \gamma / 2$, which
in turn holds whenever
$\max_{i < j} |\hat{I}_n(X_i; X_j) - I(X_i; X_j)| < \gamma / (2(p-1))$.
Hence
\[
P(\hat{T}_n \ne T^*) \;\le\;
P\!\left(\max_{i < j}|\hat{I}_n(X_i; X_j) - I(X_i; X_j)| \ge
\gamma / (2(p-1))\right).
\]
Applying a union bound over the $\binom{p}{2}$ pairs combined with
the sub-Gaussian tail \eqref{eq:subgauss} evaluated at
$t = \gamma / (2(p-1))$, and using
$2 \binom{p}{2} = p(p-1) \le p^2$, we obtain the stated bound.
Setting the right-hand side equal to $\delta$ and solving for $n$
yields
\begin{equation*}
n \;\ge\; \frac{8 \sigma^2 (p-1)^2 \bigl(2 \log p + \log(1/\delta)\bigr)}{\gamma^2}.
\end{equation*}
\qed

\subsection{Per-swap bound}
\label{app:cor2}

For each non-tree edge $e' \notin T^*$, let
$\mathrm{Cyc}_{T^*}(e')$ denote the fundamental cycle of $e'$ relative
to $T^*$, defined as the unique cycle obtained by adding $e'$ to
$T^*$. Optimality of Kruskal's algorithm on the population implies
that for every $e' \notin T^*$ and every
$e \in \mathrm{Cyc}_{T^*}(e') \cap T^*$, $I(e) \ge I(e')$. Under (A2)
the inequality is strict. We define the per-swap mutual information
gap as
\begin{equation}
\gamma_{\mathrm{swap}} \;=\; \min_{e' \notin T^*}
\min_{e \in \mathrm{Cyc}_{T^*}(e') \cap T^*} \bigl[I(e) - I(e')\bigr] > 0.
\label{eq:gamma_swap}
\end{equation}

\begin{corollary}[Per-swap finite-sample bound]
\label{cor:swap}
Under (A1)--(A2) and the per-pair sub-Gaussian tail
\eqref{eq:subgauss} with common parameter $\sigma > 0$,
\begin{equation}
P(\hat{T}_n \ne T^*) \;\le\; 2 p^3 \exp\!\left(-\frac{n
\gamma_{\mathrm{swap}}^2}{8 \sigma^2}\right).
\end{equation}
The associated sample complexity is
$n \ge 8 \sigma^2 \bigl(3 \log p + \log(2/\delta)\bigr) /
\gamma_{\mathrm{swap}}^2 = O(\sigma^2 \log p /
\gamma_{\mathrm{swap}}^2)$.
\end{corollary}

\emph{Proof.} By the matroid optimality characterisation of
maximum spanning trees, if $\hat{I}_n(e) > \hat{I}_n(e')$ for every
non-tree edge $e' \notin T^*$ and every tree edge
$e \in \mathrm{Cyc}_{T^*}(e') \cap T^*$, then $T^*$ is the unique
maximum spanning tree under $\hat{I}_n$ and Kruskal's algorithm
returns $\hat{T}_n = T^*$. Taking the contrapositive,
$\hat{T}_n \ne T^*$ implies that for at least one such pair
$(e, e')$, $\hat{I}_n(e) \le \hat{I}_n(e')$.

For a single pair $(e, e')$ with population gap
$\gamma_{e, e'} = I(e) - I(e') > 0$,
\begin{align*}
& P\bigl(\hat{I}_n(e) \le \hat{I}_n(e')\bigr) \\
& \le P\bigl(|\hat{I}_n(e) - I(e)| \ge \gamma_{e, e'}/2\bigr)
   + P\bigl(|\hat{I}_n(e') - I(e')| \ge \gamma_{e, e'}/2\bigr) \\
& \le 4 \exp\!\left(-\frac{n \gamma_{e, e'}^2}{8 \sigma^2}\right).
\end{align*}
Each non-tree edge $e'$ has a fundamental cycle of length at most
$p$, containing at most $p - 1$ tree edges (the cycle consists of
$e'$ together with the unique tree path between its endpoints). The number of
non-tree edges is at most $\binom{p}{2} - (p-1) \le p^2/2$. A union
bound over all relevant pairs, using
$\gamma_{e, e'} \ge \gamma_{\mathrm{swap}}$, yields
\begin{equation*}
P(\hat{T}_n \ne T^*) \le \frac{p^2}{2} \cdot (p - 1) \cdot 4
\exp\!\left(-\frac{n \gamma_{\mathrm{swap}}^2}{8 \sigma^2}\right)
\le 2 p^3 \exp\!\left(-\frac{n \gamma_{\mathrm{swap}}^2}{8 \sigma^2}\right).
\end{equation*}
\qed

\begin{remark}
For Kruskal's algorithm under (A2) uniqueness, the second-best
spanning tree differs from $T^*$ by a single edge swap, so the
global gap $\gamma$ and the per-swap gap $\gamma_{\mathrm{swap}}$
coincide at the population level. \Cref{cor:swap} is then
tighter than \cref{cor:tree} by a factor of $(p-1)^2$ in
the exponent. The bound is numerically informative for problems with
$\gamma_{\mathrm{swap}}$ bounded away from zero; we discuss the
practical reach of the bound on our random-mixed-tree generator in
\cref{sec:discussion} (Limitations).
\end{remark}

\section{Empirical variance-scale diagnostic}
\label{app:sigma}

The bounds in \cref{cor:tree,cor:swap} are
conditional finite-sample bounds: they hold whenever the pairwise
mutual information estimator satisfies the sub-Gaussian deviation
inequality $P(|\hat{I}_n - I| > t) \le 2 \exp(-n t^2 / (2 \sigma^2))$
with the stated $\sigma$. Finite-sample sub-Gaussian constants for
the KSG and mixed-pair $k$-NN estimators have not, to our knowledge,
been established in the literature. We do not prove the sub-Gaussian
condition. Instead, we report (i) a worst-case McDiarmid-type
calculation, (ii) an empirical variance-scale diagnostic, and (iii) a
multi-quantile tail check on the diagnostic. The latter two are
empirical and do not constitute a proof of sub-Gaussian behaviour.
The discrete--discrete plug-in estimator (used on Asia) is bounded
above by $\log\min(K_i, K_j)$ and so satisfies a sub-Gaussian tail
by Hoeffding's inequality with a constant depending on the support
sizes; we do not separately calibrate this case.

\paragraph{Worst-case theoretical bound (McDiarmid).} Each summand in
the KSG estimator is $\psi(n_x(i) + 1) + \psi(n_y(i) + 1)$, where
$\psi$ is the digamma function. The summand is bounded above by
$\psi(N) \le \log(N+1)$. Changing one data point
changes the $k$-nearest-neighbour structure of at most $O(k)$ other
points, so the bounded-difference constants are
$c_i = O((k \log N)/N)$. McDiarmid's inequality
$P(|\hat{I}_n - \mathbb{E}\hat{I}_n| > t) \le 2 \exp(-2 t^2 /
\sum_i c_i^2)$ then gives a sub-Gaussian tail with
$\sigma^2 = O\bigl((k \log N)^2 / 4\bigr)$. For $k = 5$ and $N = 1{,}500$, this gives $\sigma \approx
k \log(N) / 2 \approx 18$ nats, which is too loose to be
informative in practice.

\paragraph{Empirical variance-scale diagnostic.} The McDiarmid bound
above is a worst-case constant derived from the bounded-difference
property of the estimator and is too loose to be useful in practice.
A practical alternative is to estimate the asymptotic variance
scale directly. The asymptotic-variance relation
$\mathrm{Var}(\hat{I}_n) \approx \sigma^2 / n$ suggests an empirical
diagnostic. We draw $\hat{I}_n$ from 200
independent samples of size $n$, compute the sample variance, and
report $\sigma_{\mathrm{emp}} = \sqrt{n \cdot \mathrm{Var}(\hat{I}_n)}$.
We caution that this diagnostic estimates a variance scale; a bounded
variance does not, on its own, imply a sub-Gaussian tail (a
distribution can have bounded variance and polynomial rather than
Gaussian tails). The reported $\sigma_{\mathrm{emp}}$ is therefore the
$\sigma$ that one would substitute if the sub-Gaussian deviation
inequality held, not an established constant. The procedure is
repeated for several sample sizes $n \in \{200, 500, 1{,}500,
5{,}000\}$ to verify the $n^{-1/2}$ scaling.

\begin{table}[ht]
\centering
\caption{Empirical $\sigma$ for the KSG continuous--continuous
estimator on Gaussian pairs.}
\label{tab:sigma_cc}
\begin{tabular}{ccc}
\toprule
$\rho$ & True MI & Median $\sigma_{\mathrm{emp}}$ \\
\midrule
0.00 & 0.0000 & 0.34 \\
0.30 & 0.0472 & 0.63 \\
0.50 & 0.1438 & 0.77 \\
0.70 & 0.3367 & 0.94 \\
0.90 & 0.8304 & 1.10 \\
\bottomrule
\end{tabular}
\end{table}

\begin{table}[ht]
\centering
\caption{Empirical $\sigma$ for the mixed-pair $k$-NN estimator on
Gaussian-shift data with $K = 3$ discrete categories.}
\label{tab:sigma_cd}
\begin{tabular}{ccc}
\toprule
effect & True MI & Median $\sigma_{\mathrm{emp}}$ \\
\midrule
0.5 & 0.0771 & 0.76 \\
1.0 & 0.2538 & 0.89 \\
1.5 & 0.4472 & 0.85 \\
2.0 & 0.6191 & 0.82 \\
\bottomrule
\end{tabular}
\end{table}

The empirical $\sigma$ is approximately constant in $n$, which
confirms the $n^{-1/2}$ scaling of
$\sqrt{\mathrm{Var}(\hat{I}_n)}$. It increases with the underlying MI
for the continuous--continuous estimator and is approximately constant
in effect size for the mixed-pair estimator. Across all settings tested, $\sigma_{\mathrm{emp}} \le 1.2$.
Illustratively substituting the empirical variance-scale value
$\sigma_{\mathrm{emp}} = 1.2$ into the \cref{cor:swap}
sample-complexity formula gives
\begin{equation}
n \ge \frac{8 \cdot 1.2^2 \bigl(3 \log p + \log(2/\delta)\bigr)}{
\gamma_{\mathrm{swap}}^2}
= \frac{11.52 \bigl(3 \log p + \log(2/\delta)\bigr)}{\gamma_{\mathrm{swap}}^2}.
\end{equation}
For $p = 20$ and $\delta = 0.1$, this evaluates to
$n \ge 138 / \gamma_{\mathrm{swap}}^2$. At
$\gamma_{\mathrm{swap}} = 0.1$ nats, $n \ge 13{,}800$; at
$\gamma_{\mathrm{swap}} = 0.3$ nats, $n \ge 1{,}530$, which is
comparable to the $n = 1{,}500$ used in the experiments.

\paragraph{Multi-quantile tail check.} The variance-scale estimate
above does not test the tail behaviour. As a falsifiable diagnostic,
we draw 2000 estimates of $\hat{I}_n$ on Gaussian pairs at
$\rho = 0.5$ (true MI $\approx 0.144$) and $n = 1{,}500$, and compare
the empirical exceedance probabilities
$\hat{P}\bigl(|\hat{I}_n - I| > t\bigr)$ to the Gaussian bound
$2 \exp(-n t^2 / (2 \sigma_{\mathrm{emp}}^2))$ at several $t$ values.
For all tested $t$, the empirical exceedance probabilities are
bounded above by the corresponding Gaussian bound
(\cref{tab:quantile}).

\begin{table}[ht]
\centering
\caption{Multi-quantile diagnostic for the KSG estimator at
$\rho = 0.5$, $n = 1{,}500$, 2000 trials. Empirical exceedance
probability vs.\ the Gaussian bound at the empirically estimated
variance scale $\sigma_{\mathrm{emp}} = 0.74$. Rows with bound $> 1$
are vacuous (probabilities are bounded by 1 by definition) and are
omitted; the informative range begins around $t = 0.03$.}
\label{tab:quantile}
\begin{tabular}{ccc}
\toprule
$t$ & $\hat{P}(|\hat{I}_n - I| > t)$ & $2\exp(-n t^2/(2\sigma_{\mathrm{emp}}^2))$ \\
\midrule
0.030 & 0.110 & 0.590 \\
0.040 & 0.038 & 0.228 \\
0.050 & 0.009 & 0.067 \\
0.100 & 0.000 & 0.000 \\
0.200 & 0.000 & 0.000 \\
\bottomrule
\end{tabular}
\end{table}

At every informative $t$, the empirical exceedance probability is
bounded above by the Gaussian bound. This is consistent with
sub-Gaussian behaviour for this estimator on this distribution, but
it does not prove sub-Gaussian behaviour in general.

\paragraph{Scope of the diagnostic.} The empirical
$\sigma_{\mathrm{emp}}$ is computed on Gaussian-pair and Gaussian-%
shift conditional families only. We do not test the bimodal, skewed, or heavy-tailed conditional
families used in \cref{subsec:miacc}; the calibration is
therefore narrower than the experiments in
\cref{sec:experiments}.

Rigorous finite-sample concentration for the KSG estimator under
explicit smoothness assumptions has been studied in related
directions by \citet{singh2014exponential} for $k$-NN entropy
estimators, \citet{berrett2019efficient} for efficient $k$-NN
entropy under analytic densities, and \citet{han2020optimal} for
nonparametric divergence estimators. A formal derivation of $\sigma$
as a function of density-smoothness parameters along these lines
would convert our empirical diagnostic into a theorem.

\section{\texorpdfstring{Exploratory adaptive-$k$ table}{Exploratory adaptive-k table}}
\label{app:adaptive}

As an exploratory check using $10$ trials per sample size, we test
the rule $k(n) = \max\bigl(5, \lceil 150 / \sqrt{n}\rceil\bigr)$ on
the same MI-estimator sample-size sweep used in
\cref{subsec:suzuki}.

\begin{table}[ht]
\centering
\caption{Adaptive-$k$ variant against the MDL-discretisation MI
estimator \citep{suzuki2017chow}, $p = 15$, $\text{effect} = 0.6$,
10 trials per row (exploratory).}
\label{tab:adaptive}
\begin{tabular}{ccccc}
\toprule
$n$ & $k(n)$ & $F_1$ adaptive-$k$ & $F_1$ MDL discr. & paired $P$ \\
\midrule
75   & 18 & 0.471 $\pm$ 0.097 & 0.566 $\pm$ 0.109 & 0.0068 \\
100  & 15 & 0.471 $\pm$ 0.125 & 0.526 $\pm$ 0.112 & 0.2163 \\
200  & 11 & 0.557 $\pm$ 0.149 & 0.619 $\pm$ 0.112 & 0.0711 \\
500  & 7  & 0.664 $\pm$ 0.120 & 0.671 $\pm$ 0.107 & 0.8760 \\
1000 & 5  & 0.679 $\pm$ 0.140 & 0.629 $\pm$ 0.119 & 0.2259 \\
1500 & 5  & 0.664 $\pm$ 0.079 & 0.664 $\pm$ 0.096 & 1.0000 \\
\bottomrule
\end{tabular}
\end{table}

The adaptive rule narrows but does not eliminate the small-sample
advantage of the MDL-discretisation MI estimator. The trial count is
smaller than that used in the main MI-estimator comparison
(\cref{tab:suzuki_n}), so the adaptive-$k$ check is presented as
exploratory only and is not used to revise the main $n = 1{,}500$
comparison.

\section{\texorpdfstring{Proof of \cref{prop:forest}}{Proof of Proposition 2}}
\label{app:prop2}

We partition the $\binom{p}{2}$ pairs into
$\mathcal{E}_+ = \{(i,j) \in F^*\}$,
$\mathcal{E}_0 = \{(i,j): i, j \text{ lie in distinct components of } F^*\}$,
and $\mathcal{E}_c$ (the remaining pairs). By (A4),
$I \ge \delta > 0$ on $\mathcal{E}_+$ and $I = 0$ on $\mathcal{E}_0$.

\emph{Step 1.} For $(i,j) \in \mathcal{E}_+$, (A3) together with
$\tau_n \to 0$ implies that for sufficiently large $n$,
$\tau_n < \delta/2$, so $P(\hat{I}_n > \tau_n) \to 1$.

\emph{Step 2.} For $(i,j) \in \mathcal{E}_0$, (A5) gives
$P(|\hat{I}_n| > \tau_n) \to 0$, so $P(\hat{I}_n \le \tau_n) \to 1$.
Note that (A5) is strictly stronger than (A3). The condition
(A3) only gives pointwise convergence $\hat{I}_n
\overset{p}{\to} 0$ at any fixed threshold, but here $\tau_n$
itself shrinks to zero; (A5) is the joint requirement that the
estimator's convergence rate dominates the shrinkage rate of
$\tau_n$. The accompanying condition $\tau_n \sqrt{n} \to \infty$
ensures that $\tau_n$ dominates the $O_p(n^{-1/2})$ fluctuations
of the estimator.

\emph{Step 3.} Pairs in $\mathcal{E}_c$ have both endpoints in
the same component of $F^*$. On the event from Step 2, where
every cross-component edge in $\mathcal{E}_0$ has been removed by
the threshold rule, Kruskal's algorithm reduces to an
independent maximum-spanning-tree problem on each connected
component of $F^*$, because no surviving candidate edge crosses
components.

Fix one component $C$ and consider this restricted problem on
$C$:
\begin{itemize}
\item the vertex set is $C$;
\item the candidate edges are
$\{(i,j) \in \mathcal{E}_+ \cup \mathcal{E}_c : i, j \in C\}$;
\item the population edge weights are the pairwise mutual
informations $I(X_i; X_j)$ for $i, j \in C$;
\item the population target is the maximum spanning tree on $C$
under these weights, which by (A2$'$) is unique and equals the
restriction of $F^*$ to $C$.
\end{itemize}

The hypotheses of \cref{prop:main} hold on this restricted
problem: (A1) is invariant under restriction, (A3) is a pairwise
condition that still holds for every pair in $C$, and (A2$'$)
gives uniqueness of the target tree on $C$. Applying
\cref{prop:main} to the restricted problem therefore yields that
Kruskal's algorithm run on the within-component candidate edges
returns the restriction of $F^*$ to $C$ with probability tending
to 1.

Repeating this argument for every component of $F^*$ gives that
the empirical forest matches $F^*$ on $\mathcal{E}_+$ with
probability tending to 1. Any further edge from $\mathcal{E}_c$
considered later by Kruskal would close a cycle within an
already-built component, and is therefore rejected.

\emph{Step 4.} The intersection of the events in Steps~1 to 3 has
probability tending to 1. On this event, Kruskal's algorithm with
threshold $\tau_n$ returns exactly $F^*$ by (A2$'$). \qed

%%%%%%%%%%%%%%%%%%%%%%%%%%%%%%%%%%%%%%%%%%%%%%%%%%%%%%%%%%%%%%%%%%%%%%%%
% Bibliography. The references are stored in references.bib in the
% same folder. For Behaviormetrika submission, you may need to swap
% \bibliographystyle{plainnat} for the journal's required style
% (e.g., spbasic).
\bibliographystyle{plainnat}
\bibliography{references}

\end{document}
